\begin{lemma}{Gauß'scher Min--Max--Satz für Gordon'sche Paare (endlich)}
             {GordonschePaareUndDerGaussscheMinMaxSatzEndlich}
Seien $\emptyset \neq T \subseteq S^{m-1}$ und
$\emptyset \neq V \subseteq \R^N\setminus\meng 0$ beschränkt.\\
Seien $c \in \R_{>0}$ mit
$\forall\, v \in V: \norm v _{2, N} \le c$
und $(X, Y)$ das Gordon'sche Paar zum Tupel $(\gamma, \delta, g, c)$.
Schließlich seien $\mct := \Pot_{ne}(T)$ und $\mcv := \Pot_{ne}(V)$.\\
Dann gilt:
\begin{enumerate}[(1)]
	\item \label{GordonschePaareUndDerGaussscheMinMaxSatzEndlich1}
	$\forall\, A \in \mct\, \forall\, B \in \mcv:
	\E\left(\Inf{t \in A}\Sup{v \in B}X(t, v) \right)
	\le \E\left(\Inf{t \in A}\Sup{v \in B}Y(t, v) \right)$.
	\item \label{GordonschePaareUndDerGaussscheMinMaxSatzEndlich2}
	$\forall\, A \in \mct\, \forall\, B \in \mcv:
	\E\left(\Sup{t \in A}\Sup{v \in B}Y(t, v) \right)
	\le \E\left(\Sup{t \in A}\Sup{v \in B}X(t, v) \right)$.
\end{enumerate}
\end{lemma}
\begin{proof}
Seien $A \in \mct$ und $B \in \mcv$.\\
Da $A$ und $B$ endlich und nichtleer sind, existieren $r, s \in \N$ und
$\alpha : \haken r \rightarrow A$ bijektiv und $\beta: \haken s \rightarrow B$
bijektiv.\\
Seien $i, j \in \haken r$ und $k, l \in \haken s$.\\
Nach \refclaim{IntegrierbarkeitUndVerteilungVonLinearkombinationen}
{IntegrierbarkeitUndVerteilungVonLinearkombinationen1} sind 
$X(\alpha(i), \beta(j))$ und $Y(\alpha(i), \beta(j))$ quadratintegrierbar und
damit
\begin{align*}
X'& : \haken r \times \haken s \rightarrow L_2(\Omega, \mfa, \Prim, \R),\ 
(a, b) \mapsto X(\alpha(a), \beta(b))\\
Y'& : \haken r \times \haken s \rightarrow L_2(\Omega, \mfa, \Prim, \R),\ 
(a, b) \mapsto Y(\alpha(a), \beta(b))
\end{align*}
wohldefiniert.\\
Da $\alpha$ und $\beta$ bijektiv, sowie $A$ und $B$ endlich sind, gilt
offenbar:
\begin{align}\inf_{t \in A}\sup_{v \in B}X(t, v) =
  \min_{a \in \haken r}\max_{b \in \haken s}X(\alpha(a), \beta(b))
  = \min_{a \in \haken r}\max_{b \in \haken s}X'(a, b) \label{gpugmmseeq:1}
\end{align}
und analog:
\begin{align}
\inf_{t \in A}\sup_{v\in B}Y(t, v) 
= \min_{a \in \haken r}\max_{b \in \haken s}Y'(a, b)\text,\label{gpugmmseeq:2}
\end{align}
sowie
\begin{align}\sup_{t \in A}\sup_{v\in B}X(t, v) 
&=\max_{a\in \haken r}\max_{b \in \haken s}X'(a, b)\text,\label{gpugmmseeq:3}\\
\sup_{t \in A}\sup_{v\in B}Y(t, v) 
&=\max_{a \in \haken r}\max_{b \in \haken s}Y'(a, b)\text.\label{gpugmmseeq:4}
\end{align}
Wegen $\alpha(i) \in S^{m-1}$ ist $\alpha(i) \neq 0$ und wegen
$\beta(i) \in V \subseteq \R^N\setminus\meng 0$ ist $\beta(i) \neq 0$.\\
Somit ist nach
\refclaim{IntegrierbarkeitUndVerteilungVonLinearkombinationen}
{IntegrierbarkeitUndVerteilungVonLinearkombinationen2}: 
$X(\alpha(i), \beta(j))\neq 0$ und nach
\refclaim{IntegrierbarkeitUndVerteilungVonLinearkombinationen}
{IntegrierbarkeitUndVerteilungVonLinearkombinationen3}:
$Y(\alpha(i), \beta(j)) \neq 0$.\\
Mit $\norm{\alpha(i)}_{2,m}=1$ gilt nach
\refclaim{IntegrierbarkeitUndVerteilungVonLinearkombinationen}
{IntegrierbarkeitUndVerteilungVonLinearkombinationen4}
\[X'(i,j)=X(\alpha(i), \beta(j)) \ \text{ ist }\ 
N(0, c^2 + \norm {\beta(j)}_{2, N}^2)\text{--verteilt}\] und 
nach \refclaim{IntegrierbarkeitUndVerteilungVonLinearkombinationen}
{IntegrierbarkeitUndVerteilungVonLinearkombinationen5}
\[Y'(i, j)=Y(\alpha(i), \beta(j)) \ \text{ ist } \ 
N(0, \norm{\beta(i)}_{2, N}^2)\text{--verteilt.}\]
Weiter gilt ebenfalls wegen $\alpha(i) \in S^{m-1}$:
\begin{align*}
&\, \norm{X'(i, k) - X'(i, l)}_2  = \norm{X(\alpha(i), \beta(k)) 
- X(\alpha(i), \beta(l))}_2\\ \overset{\refclaim{AbstaendeInGordonschenPaaren}
{AbstaendeInGordonschenPaaren4}} = &\, \norm{Y(\alpha(i), \beta(k)) 
- Y(\alpha(i), \beta(l))}_2 = \norm{Y'(i, k) - Y'(i, l)}_2\text.
\end{align*}
Es ist auch $\alpha(j) \in S^{m-1}$, sowie 
$\beta(k), \beta(l)\in V$ und damit nach Voraussetzung 
$\norm{\beta(k)}_{2,N} \le c$, $\norm{\beta(l)}_{2,N} \le c$, also gilt:
\begin{align*}
&\, \norm{Y'(i, k) - Y'(j, l)}_2 
= \norm{Y(\alpha(i), \beta(k)) - Y(\alpha(j), \beta(l))}_2\\
\overset{\refclaim{AbstaendeInGordonschenPaaren}
{AbstaendeInGordonschenPaaren5}}\le &\,
\norm{X(\alpha(i), \beta(k)) - X(\alpha(j), \beta(l))}_2
= \norm{X'(i, k) - X'(j, l)}_2\text.
\end{align*}
Damit erfüllen $X', Y'$ alle Voraussetzungen von
\ref{GaussscherMinMaxSatzGleichheitInDerErstenBedingung}, sodass gilt:
\begin{align*}
\E\left(\inf_{t \in A}\sup_{v \in B}X(t, v) \right)
&\overset{\eqref{gpugmmseeq:1}}=
\E\left(\min_{a \in \haken r}\max_{b \in \haken s}X'(a, b)\right)
\overset{\refclaim{GaussscherMinMaxSatzGleichheitInDerErstenBedingung}
{GaussscherMinMaxSatzGleichheitInDerErstenBedingung1}}\le
\E\left(\min_{a \in \haken r}\max_{b \in \haken s}Y'(a, b)\right)\\
& \overset{\eqref{gpugmmseeq:2}}= 
\E\left(\inf_{t \in A}\sup_{v \in B}Y(t, v) \right)
\end{align*}
und
\begin{align*}
\E\left(\sup_{t \in A}\sup_{v \in B}X(t, v) \right)
&\overset{\eqref{gpugmmseeq:3}}=
\E\left(\max_{a \in \haken r}\max_{b \in \haken s}X'(a, b)\right)\\
&\overset{\refclaim{GaussscherMinMaxSatzGleichheitInDerErstenBedingung}
{GaussscherMinMaxSatzGleichheitInDerErstenBedingung2}}\le
\E\left(\max_{a \in \haken r}\max_{b \in \haken s}Y'(a, b)\right)
\overset{\eqref{gpugmmseeq:4}}=
\E\left(\sup_{t \in A}\sup_{v \in B}Y(t, v) \right)\text.
\end{align*}
\end{proof}